\begin{figure}[H]
\centering
\begin{tikzpicture}[
  box/.style={draw=linkcolor,rounded corners=2pt,align=center,
    font=\small,text width=4.8cm,inner sep=7pt},
  wide/.style={box,text width=11.1cm},
  flow/.style={-{Latex[length=2mm]},draw=linkcolor,thick},
  node distance=9mm and 7mm]
  \node[wide] (smooth) {Boundedness and cocompactness\\
    Smoothness of $U$; proper intrinsic distance; disc continuity};
  \node[wide,below=of smooth] (boundary) {Semialgebraic boundary preparation\\
    Joint normal-offset limits at a point with independent peaks};
  \node[box,anchor=north east] (supports)
    at ([xshift=-3.5mm,yshift=-9mm]boundary.south)
    {Independent peaks\\Control every scaled normal coordinate};
  \node[box,anchor=north west] (discs)
    at ([xshift=3.5mm,yshift=-9mm]boundary.south)
    {Regularized analytic discs\\Interior balls prevent collapse};
  \node[wide,below=35mm of boundary] (scaling)
    {Two-sided weighted scaling\\
    Interior inclusion $+$ excluded points $\Longrightarrow$ global equivalence};
  \node[wide,below=of scaling] (vey)
    {Quadratic model with a bounded realization and compact quotient\\
    Vey's theorem $\Longrightarrow$ bounded symmetric domain};
  \draw[flow] (smooth) -- (boundary);
  \draw[flow] (boundary.south) -- ++(0,-3mm) -| (supports.north);
  \draw[flow] (boundary.south) -- ++(0,-3mm) -| (discs.north);
  \draw[flow] (supports.south) -- ++(0,-3mm) -| (scaling.north);
  \draw[flow] (discs.south) -- ++(0,-3mm) -| (scaling.north);
  \draw[flow] (scaling) -- (vey);
\end{tikzpicture}
\caption{The two boundary inputs have different roles. Peaks bound the
forward scaling maps, while discs ensure a nonempty normal interior.
The two-sided convergence then identifies the global quadratic model.
}
\label{fig:proof}
\end{figure}
